\chapter{Données}

Dans le cadre de ce mémoire, les données textuelles exploitées reposent sur trois ressources principales : 
le projet ParCoLab, l’Atlas sonore des langues régionales du LIMSI, et une traduction de la parabole du 
fils prodigue. Ces corpus ont été choisis pour leur richesse linguistique, correspondant aux attentes de
notre projet. En effet, ces 3 projets s'intéressent aux langues régionales de France particulièrement alsacien,
corse et occitan.

Le corpus ParCoLab (ParCoLab – Multilingual Parallel Corpus) \citep{stosic_2024} constitue la base la plus 
volumineuse. ParCoLab est un projet parallèle multilingue porté par le laboratoire CLLE 
(Université de Toulouse) et soutenu par la Délégation générale à la langue 
française et aux langues de France (DGLFLF). Celui-ci vise à 
constituer et diffuser un corpus de textes alignés dans différentes
langues, dont des langues régionales de France centrales du projet.
Ces textes sont alignés au niveau des phrases, ce qui permet d’effectuer des analyses comparatives fines. 

L’Atlas sonore des langues régionales du LIMSI (Laboratoire LIMSI, CNRS) \citep{boulademareuil_2017} est 
une ressource patrimoniale remarquable. Ce projet, consiste en une 
carte interactive de France sur laquelle sont disséminés plus de 200 commune où 
des locuteurs ont enregistré des versions d’un même texte : la fable d’Ésope « La bise et le soleil ».
Ces enregistrements, souvent retranscrits dans des graphies adaptées, permettent d’étudier des dialectes 
des diverses langues régionales françaises, et d’analyser la variation phonétique, lexicale ou morphologique 
dans un contexte de terrain. Grâce à cette ressource, nous avons pu intégrer des textes en occitan très variés 
phonétiquement, enrichissant le corpus avec des variations dialectiques.

Enfin, j’utilise une traduction de la parabole du fils prodigue (également connue sous le nom de 
“parabole de l’enfant prodigue”), réalisée dans différentes variétés régionales. Cette parabole est 
un classique des enquêtes dialectologiques : elle a été traduite dans de nombreuses parlers régionaux, 
notamment via les enquêtes menées par les Coquebert de Montbret au XIXe siècle. Sur le plan linguistique, 
la version occitane de cette parabole (ainsi que des variantes en corse ou d’autres langues) offre un point 
d’ancrage commun — un texte “pivot” — qui facilite la comparaison entre dialectes et la détection de 
structures grammaticales, lexicales ou morphologiques typiques de chaque variété.

\section{Analyse de données}

\subsection{Données parallèles}

\begin{table}[h!]
\centering
\begin{tabular}{lccc}
\hline
\textbf{Données} & \textbf{Atlas LIMSI} & \textbf{Parcolab} & \textbf{Parabole du fils prodigue} \\
\hline
Tokens & 7\,922 & 35\,371 & 11\,729 \\
Phrases & 305 & 1\,485 & 537 \\
Lemmes fr & 57 & 2\,791 & 191 \\
Textes fr & 1 & 9 & 1 \\
Textes co & 6 & 9 & -- \\
Textes gsw & 7 & 9 & -- \\
Textes oc & 47 & 9 & 21 \\
\hline
\end{tabular}
\caption{Données statistiques des trois corpus parallèles}
\label{table_donnees_para}
\end{table}

Les données du tableau \ref{table_donnees_para} sont issus d'un programme python fait par nos soins. 
Ce dernier utilise la bibliothèque Spacy de manière à trouver le nombre
de phrases, de tokens (hors poncutations, avec apostrophes et contractions séparés en deux tokens), 
ainsi que le nombre de lemmes français. Aucun modèle spacy n'existe pour l'alsacien, le corse et l'occitan,
ce qui rend impossible une lemmatisation. Cette difficulté est exacerbée car ces langues sont principalement oral
et ne possède pas de réelles règles ortographiques.
On retrouve dans notre corpus une majorité de textes en occitan (77), mais aussi 15 textes en corse et 16 
textes en alsacien. Le fait que ParCoLab soit le corpus le plus textuellement riche (dû à la taille des textes 
présents dans celui-ci) explique qu’il présente également le nombre de tokens le plus élevé.

\begin{table}[h!]
\centering
\begin{tabular}{lcccc}
\hline
\textbf{Données} & \textbf{Français} & \textbf{Alsacien} & \textbf{Corse} & \textbf{occitan} \\
\hline
Fichiers & 11 & 16 & 15 & 77 \\
Tokens & 10\,070 & 9\,677 & 9\,601 & 25\,674 \\
Phrases & 363 & 423 & 418 & 1\,123 \\
Lemmes & 3039 & -- & -- & -- \\
\hline
\end{tabular}
\caption{Données statistiques du corpus parallèles par langues}
\label{table_donnees_par_langues}
\end{table}

On peux remarquer avec ce second tableau (\ref{table_donnees_par_langues}) comparant la répartitions des fichiers, tokens et phrases
précedemments établis en fonction des langues, que la langue avec la plus grande variété dialectique et
l'occitan. Cette langue possédant la vaste majorité des tokens et un nombre de fichiers bien plus haut
(justement dû au nombre de ces variétés présentes dans l'atlas LIMSI). Le corse et l'alsacien sont
beaucoup plus proche statistiquement et sont même très similaire. Le français nous présente un nombre de
lemmes qui devrait, en théorie, être similaire peu importe la langue utilisé (les textes étant des traductions). 
Cependant la lemmatisation des autres langues reste pour l'instant impossible dû au manque de modèle Spacy
pour celles-ci.

\subsection{Données comparables}

\begin{table}[h!]
\centering
\begin{tabular}{lcccc}
\hline
\textbf{Données} & \textbf{Alsacien} & \textbf{Corse} & \textbf{Francais} & \textbf{Occitan} \\
\hline
Fichiers & 2\,399 & 22 & 3\,487 & 29\,49 \\
Tokens & 528\,337 & 9004 & 5254875 & 380\,374 \\
Phrases & 33\,814 & 437 & 235\,643 & 17\,239 \\
Lemmes & -- & -- & 1\,519\,413 & -- \\
\hline
\end{tabular}
\caption{Données statistiques du corpus comparable Wikimedia}
\label{table_donnees_comp}
\end{table}

Le tableau \ref{table_donnees_comp} représente les données sur lesquelles le modèle entraîné va-t-être validé. 
Celles-ci ont étaient obtenue via Wikimedia et les liens interlangues. 
Pour ce faire nous avons utilisé le programme de \citet{tsolakis_2024} que nous avons modifié afin d'avoir 
les pages Wikimedia en français, alsacien, occitan et corse. On remarque dans notre cas que le nombre de
tokens et bien plus conséquent, surtout dans le cas de l'alsacien qui possède une communauté très active
et inclue les dialectes mulhousien et strasbourgeois qui apporte à la variété dialectique et ajoute des
tests possibles sur les graphies dialectales.